\section{Two raw-kernel realizations and conventional tasks}\label{sec:realizations}
These verifications are consequences of the general theorems, never their premises. Their kernel mechanisms were established in the preceding article\cite{prior}; we include complete proofs and specify ordinary prediction tasks.

\begin{lemma}[Positive-instrument coupling]\label{lem:coupling}
Let states be probability vectors in $\ell^1$ or finite-dimensional density matrices in trace norm. Let $C$ be a positive trace-preserving affine preparation channel with contraction coefficient $\gamma$. Let $A_y$ be a positive linear instrument density with $\int A_y\,\zeta(dy)$ trace preserving. For $q=C(s)$ put $u_y=A_y(q)$, $r_y=\tr u_y$, and $F_y=u_y/r_y$ on positive-mass reports. Then
\begin{equation}\label{eq:coupling}
 W_1(\law F_Y,\law F'_{Y'})\le3\gamma\|s-s'\|,
 \qquad\TV(\law Y,\law Y')\le\frac{\gamma}{2}\|s-s'\|.
\end{equation}
\end{lemma}
\begin{proof}
The positive and negative parts of a signed vector or Hermitian matrix $v$ give
\[
 \int\|A_y(v)\|\,\zeta(dy)\le\|v\|,\qquad
 \int|\tr A_y(v)|\,\zeta(dy)\le\|v\|.
\]
Couple common report mass $\min(r_y,r'_y)$ at the same $y$. There
\[
 \min(r_y,r'_y)\|F_y-F'_y\|\le\|u_y-u'_y\|+|r_y-r'_y|.
\]
The state space has diameter at most two, so the remaining report mass costs at most $\int|r_y-r'_y|\,d\zeta$. Summing proves the Wasserstein bound; the second integral directly gives the TV bound. These are sufficient coefficients, not claimed optimal ones.
\end{proof}

\subsection{Gaussian hidden-state sensing}
There are $d+1$ hidden states, a fixed finite action family, and known stochastic matrices $T_t^a$ with every entry at least $\tau>0$. A call applies $T_t^a$ and then reports
\[
 Y\mid X'=j,a\sim N(\mu_j^a,\Sigma_a),\qquad Y\in\R^{m_{\rm obs}},\quad m_{\rm obs}\ge d.
\]
The parameter set is compact: the means are uniformly bounded, the eigenvalues of $\Sigma_a$ are uniformly bounded above and below, and the least eigenvalue of
$B_a\Sigma_aB_a^T$ is uniformly positive, where $B_a$ has rows $(\mu_j^a-\mu_0^a)^T\Sigma_a^{-1}$. At each phase at least one legal $T_t^a$ is invertible. Put
$\gamma_t=\max_a\frac12\max_{i,j}\sum_k|T_t^a(i,k)-T_t^a(j,k)|$.

\begin{theorem}[Categorical filtering]\label{thm:gaussian}
If $\beta_t=3\gamma_t$ obeys \eqref{eq:regular_stability}, all conclusions of Theorem~\ref{thm:regular} hold for the conventional loss
\[
 L(X_n,u)=\|e_{X_n}-u\|_2^2,\qquad u\in\Delta_d,
\]
with predictive carrier $\Delta_d$. The full-history baseline is the usual optimal categorical filtering Bayes risk in this controlled experiment. Constants are uniform in $n,M$ on the stated parameter class. Every Gaussian report, including its tails, is a paid acquisition.
\end{theorem}
\begin{proof}
For incoming $p$ set $q=pT_t^a$; then $\tau\le q_j\le1-d\tau$. With the Gaussian emission densities $f_j$, the raw joint kernel gives
\[
 r(q,y)=\sum_jq_j f_j(y),\qquad u_j=F_j(p,a,y)=q_j f_j(y)/r(q,y).
\]
Unnormalized finite continuation probabilities are linear in $p$, proving the instrument mixture identities on the whole simplex. All report densities are positive on $\R^{m_{\rm obs}}$, and therefore equivalent to one fixed nonsingular Gaussian probability reference for each action, at every phase.

The actual posterior law has uniformly bounded density even at the simplex boundary. Indeed its logits satisfy
\[
 z_j=\log(u_j/u_0)=(B_ay)_j+(c_a)_j+\log(q_j/q_0),\qquad
 (c_a)_j=-\tfrac12(\mu_j^T\Sigma_a^{-1}\mu_j-\mu_0^T\Sigma_a^{-1}\mu_0).
\]
Their law is a finite mixture of nonsingular $d$-dimensional Gaussians, with uniform bounds on means and covariance eigenvalues. Hence its density obeys
$h_q(z)\le C\exp(-c|z|^2+C|z|)$. The softmax map has Jacobian $\prod_{j=0}^du_j$, and
\[
 (\textstyle\prod_j u_j)^{-1}
 =(1+\sum_{j=1}^d e^{z_j})^{d+1}e^{-\sum_j z_j}
 \le C_d e^{C_d|z|}.
\]
Consequently $h_q(\operatorname{logit}u)/\prod_j u_j$ is bounded uniformly on the entire simplex interior. The boundary has zero acquired mass; no clipping or conditioning was used.

The row coefficient gives $\|pT-p'T\|_1\le\gamma_t\|p-p'\|_1$: decompose the difference into equal positive and negative masses and bound the distance of two row mixtures. Lemma~\ref{lem:coupling} proves report continuity and the required Wasserstein estimate. This also gives weak continuity of the acquired kernel.

For predictive separation, distinct Gaussian means give linearly independent emission densities: divide a purported relation by the common quadratic Gaussian factor and restrict the resulting exponentials to a line with distinct slopes. Successive limits force all coefficients to vanish. The next-report law thus identifies $q$, and an invertible legal $T_t^a$ identifies $p$. Conversely $p$ determines all future tests, so this is the predictive quotient rather than a redundant parameterization.

Finally $\ell(p,u)=1-2p\cdot u+\|u\|_2^2$ is affine in $p$ and has Bayes risk $g(p)=1-\|p\|_2^2$. It is Lipschitz and strongly concave on the affine simplex. Apply Theorem~\ref{thm:regular}. The proof optimizes over the full declared adaptive policy class; it does not infer a strict feedback advantage merely from the existence of several sensors.
\end{proof}
For instance, two-state invertible transition matrices with row coefficients repeating $0.8,0.8,0.02$ have individual certified coefficients as large as $2.4$, but every three-step product is $0.3456<1$. Weak mixing by itself is not being substituted for this explicit block condition.

\subsection{Noncommuting finite-dimensional instruments}
Let $D\ge2$, $d=D^2-1$, and take an orthonormal real basis $H_1,\ldots,H_d$ of traceless Hermitian matrices. Choose $s>0$ with $s\sum_i\|H_i\|_{\rm op}<1$. With uniform probability measure on the report cube $[-1,1]^d$, put
\[
 E_y=I+s\sum_i y_iH_i,\qquad
 \mathcal I_{t,y}^a(\rho)=E_y^{1/2}\mathcal C_t^a(\rho)E_y^{1/2},
\]
where
\[
 \mathcal C_t^a(\rho)=(1-\gamma_t^a)\tau_t^a+
                    \gamma_t^a U_a\rho U_a^*,\qquad
 0<\gamma_-\le\gamma_t^a\le\gamma_+<1,\quad \tau_t^a\succeq\tau I.
\]
The affine channel extends linearly by multiplying its first term by $\tr\rho$. It is a paid mixture of preparation and a unitary channel. Since $\int E_y\,2^{-d}dy=I$, the instrument is positive and normalized. No common eigenbasis of the effects is assumed.

\begin{theorem}[Prediction of an informationally complete measurement]\label{thm:quantum}
If $\beta_t=3\max_a\gamma_t^a$ satisfies the block condition, Theorem~\ref{thm:regular} holds on the density-matrix carrier for Brier prediction of the outcome of a fixed finite informationally complete POVM. Affine volume means Lebesgue volume in the displayed orthonormal traceless chart, and the Wasserstein metric uses trace norm. The constants need not be uniform in $D$. The terminal measurement is paid when included in the physical acquisition count.
\end{theorem}
\begin{proof}
Let $q=\mathcal C_t^a(\rho)$, $r_y=\tr(E_yq)$ and
$\sigma_y=E_y^{1/2}qE_y^{1/2}/r_y$. Uniform positive lower bounds on $q,E_y$ give equivalent report supports on the whole cube. Linearity of the unnormalized quantum instruments proves mixture closure. The coupling lemma supplies the Wasserstein coefficient $3\gamma_t^a$ and TV continuity.

For the density calculation, at positive definite $q,\sigma$ define
\[
 A=q^{-1/2}(q^{1/2}\sigma q^{1/2})^{1/2}q^{-1/2}.
\]
This is the unique positive definite solution of $AqA=\sigma$. The inverse of the normalized congruence map on the trace-$D$ positive hyperplane is
\begin{equation}\label{eq:quantum_inverse}
 E=\frac{DA^2}{\tr A^2}.
\end{equation}
Since $A$ is positive definite, $E^{1/2}=\sqrt{D/\tr A^2}\,A$, proving the substitution in both directions. The map and its inverse are smooth; their differentials on the trace hyperplanes are inverse. The affine map $y\mapsto E_y$ has full rank. The permitted $q$ and closed cube lie in a compact subset of the positive domains, so the absolute Jacobian of $y\mapsto\sigma_y$ has a positive uniform lower bound. The bounded report density and one-to-one change of variables prove a uniform acquired-density upper. Cube boundaries are null, and no multiplicity factor is hidden.

The report density $1+s\sum_i y_i\tr(H_iq)$ identifies all traceless coordinates of $q$; the affine channel is injective since $\gamma_t^a>0$. Thus the density matrix is exactly the predictive quotient for the declared tests. It is derived from the raw instruments, not supplied to the finite observer as a register.

For a concrete fixed POVM take
$G_{i,\pm}=(I\pm eH_i)/(2d)$ with
$0<e<1/\max_i\|H_i\|_{\rm op}$. These effects sum to $I$. Its outcome probabilities $P_{i,\pm}(\rho)=\tr(G_{i,\pm}\rho)$ form an injective affine map. Under ordinary Brier loss the Bayes risk is $1-\|P(\rho)\|_2^2$. Its strong-concavity constant is twice the squared least singular value of the linear part of $P$ on the traceless chart. This is a probability-forecast task for an actual fixed measurement, not an assertion that every discrimination loss has full-dimensional curvature. Theorem~\ref{thm:regular} now applies.
\end{proof}
As a finite-stratum extension, replace a fixed fraction $\vartheta<1$ of calls by a flagged preparation of one of finitely many pure states and use the displayed continuous instrument otherwise. A subsequent positive channel sends pure states back into the interior. The type is visible in the current report; the continuous terminal submeasure has its actual mass $1-\vartheta$. Provided the corresponding coupling block bound holds, Proposition~\ref{prop:strata} applies with its unnormalized mass factor. This verifies genuine rank changes without asserting that arbitrary or unobserved quantum strata satisfy that proposition.
